\chapter{Using the public workflow}
\label{sec:current-guide}

\section{Specify the tuning problem before the search}

A model parameter is fitted within one trial, such as logistic coefficients
or neural weights. A hyperparameter is selected between trials, such as a
penalty or learning rate. Architecture coordinates change the model itself,
such as width and depth. For each candidate configuration, the objective fits
using only training identities, scores the declared tuning identities and
returns a measured loss. The searcher receives the complete tuning mean.
After selection is frozen, fresh fits use the declared final folds and seeds.

Choose stable model, data, objective and protocol identities. Declare objective
direction and physical units; define train/score membership for every fold;
reserve final work separately. A known data generator supplies a mathematical
reference but does not make finite-sample hyperparameter choice analytically
obvious. The response example distinguishes those facts explicitly.

\section{The caller objective contract}

An endpoint callable accepts \texttt{config}, \texttt{seed},
\texttt{evaluation} and \texttt{trial\_id}. The execution context identifies
the role, fold and declaration. The adapter owns data loading, preprocessing,
training, artifact serialization and reload. The return shape is:

\begin{verbatim}
def objective(config, seed, evaluation, trial_id):
    train = load_rows(evaluation.fold.train_ids)
    score = load_rows(evaluation.fold.score_ids)
    model = fit(config, train, seed)
    loss = evaluate_loss(model, score)
    model_path = save_once(trial_id, model)
    return {
        "value": loss,
        "artifacts": {"model": model_path},
        "metadata": {"resource": {
            "work": measured_work, "unit": "training_steps"
        }},
    }
\end{verbatim}

This fragment is a contract template: \texttt{load\_rows}, \texttt{fit},
\texttt{evaluate\_loss}, \texttt{save\_once} and the meter are supplied by
the caller. It is not a runnable model. The installed quick starts and complete
application companion source provide executable implementations. Writes must
be idempotent under the trial identity: an interruption may repeat an
uncommitted call. Do not overwrite an incompatible existing artifact.

The public wiring is illustrated by the following executable example. It uses
the installed fitted-readout caller model; replacing that callable with a new
model is the integration point. The command-line examples in the preceding
chapter are the shorter route for first use.

\begin{verbatim}
from contextlib import closing
from pathlib import Path
from hponas.quickstart import make_study
from hponas.store import Store

directory = Path("/tmp/manual-caller-example")
directory.mkdir(exist_ok=False)
with closing(Store(directory / "study.sqlite")) as store:
    study = make_study("hinge", directory, store)
    result = study.run()
    study.export_report(directory / "report.json")
    print(result.best_config.values, result.best_value)
    print(result.evaluation["outcome"])

    resumed = make_study("hinge", directory, store)
    resumed.load_checkpoint()
    same_result = resumed.run()
    assert same_result.best_value == result.best_value
\end{verbatim}

For a custom model, first construct a \texttt{StudyDeclaration}.

Its objectives use \texttt{EvaluationObjective}; its folds use
\texttt{EvaluationFold}. Pass the searcher, executor and callable to
\texttt{Study}, together with that declaration, Store and budget.
The full tested wiring appears in
\path{hpo-survey/examples/application_examples.py}, functions
\texttt{Objective} and \texttt{make\_study}. Architecture coordinates are
explicitly supplied as \texttt{('width', 'depth')} for the neural example.
The objective must actually construct the corresponding architecture; naming
a coordinate cannot perform that work automatically.

\section{Evaluation and leakage}

The default declaration seeds are $(0,1)$ for tuning and $(1000,1001)$ for
final evaluation. Defaults provide repetitions, not precision guarantees.
Every candidate runs all declared fold/seed combinations; means weight each
combination equally even when fold sizes differ. A failed or cancelled
repetition invalidates its candidate rather than disappearing from the mean.
Vector constraints use the worst repetition. Final data cannot alter the
frozen selection or train the searcher.

\texttt{rows} folds reject train/score identity overlap. \texttt{grouped}
folds additionally separate entities. \texttt{temporal} folds declare
information availability through \texttt{train\_until}, \texttt{score\_from},
\texttt{score\_until} and a nonnegative gap. Final scoring identities cannot
train another final fold. These checks cannot inspect arbitrary undeclared
access inside caller code; the data adapter remains a trust boundary.
To estimate the performance of the entire tuning procedure, use an appropriate
outer design across separate studies. A final test of one frozen winner is not
automatically nested validation of the selection procedure.

Without a declaration, a nonempty persisted \texttt{legacy\_reason} is
mandatory and the report says \texttt{legacy\_unqualified}. A deterministic
fixture can explicitly opt out of final evaluation using
\texttt{opt\_out\_reason} and empty final folds/seeds. Neither route creates
protected-final evidence. Old snapshots cannot acquire a fabricated reason
through recovery; changed protocols require a new study or an explicit
retained migration.

\section{Configuration count, work and time}

\texttt{StudyBudget} separates configurations from calls. Folds, seeds,
promotions, final fits and retries add work. Each call reserves a declared
maximum; measured incremental work replaces that reservation on return.
Uncommitted lost work consumes the full reservation, and a retry consumes a
new one. Wrong units, nonfinite meters and exceeded bounds fail the measurement
and stop further dispatch in that role.

\begin{verbatim}
from hponas import StudyBudget
resources = StudyBudget(
    configurations=8, work_unit="training_steps",
    tuning_work=160, final_work=20,
    max_work_per_call=10, worker_capacity=1,
    tuning_seconds=60, final_seconds=20,
    full_run_work=10,
)
\end{verbatim}

This envelope can cover eight configurations, two tuning seeds and two final
seeds with ten units per call; it provides no retry headroom at that maximum.
The final reserve must cover one scalar winner or the maximum vector front.
Time is summed measured call duration, not parallel wall time. Endpoint caps
can overshoot by the final call; unknown elapsed work after a crash is not
invented. This accounting is not a hard sandbox for arbitrary objective code.
Use external process limits when needed, as the DOC-01 runner does.

\section{Streaming, multiple objectives and populations}

Repeated declared streaming \emph{is supported}. Set
\texttt{streaming\_steps} to include every scheduler endpoint and any desired
intermediate barrier. The callable accepts \texttt{reporter},
\texttt{checkpoint} and \texttt{resume\_from}; each segment advances one
repetition, reports its endpoint and returns a JSON checkpoint. Only complete
cohorts reach the scheduler. A pause retains each repetition's checkpoint;
final fits restart on final roles and their barriers never feed scheduling.
Coarse barriers restrict pruning opportunities.

For vector optimization, use \texttt{MultiObjectiveStudy} with an explicit
\texttt{ObjectiveSpec}; return named objective/constraint mappings and work.
The result is a feasible front, possibly empty. Do not conceal failed feasibility
inside a small scalar training loss. Population adapters additionally implement
model/optimizer state transfer and compatibility checks; saved lineage and
schedules are first-class outputs. Finite local population support does not
require the proposal's former eight-worker boundary or prove large-scale RL.

\section{Recover and inspect}

Retain matching code, dependency versions, declaration, searcher configuration,
database/checkpoint and caller artifacts. Construct fresh searcher and executor
objects, attach the same Store and call \texttt{load\_checkpoint()} before
\texttt{run()}. Store snapshots are authoritative when both Store and file
projections exist. Evaluation checkpoints use schema 4; resource-declared
checkpoints use schema 5. A report cannot substitute for these snapshots.

\texttt{study.report()} returns a detached report;
\texttt{study.export\_report(path)} and the Store export interface refuse an
existing destination. Reports retain complete raw roles, means, sample spread,
counts, resource uncertainty, artifacts and tuning trajectories. Nonfinite
values become JSON null with explicit classifications, not silently valid
numbers. \texttt{reported} indicates that a workflow ended; inspect
\texttt{outcome} for completed, failed, infeasible or no-selection status.

Rung diagnostics compare observed completed full-fidelity tuning pairs and
retain counts and censoring reasons. Protected final data is excluded. A rank
correlation does not certify that pruning is safe on a different task.

For Ray Core with \texttt{StudyBudget}, set \texttt{max\_retries=0}; hidden
executor retries would bypass reservation accounting. Local recovery,
PostgreSQL and service migrations have their own contracts in
\path{docs/OPERATIONS.md}, \path{docs/POSTGRES.md} and
\path{docs/REPORTS_AND_ADAPTERS.md}. Application roles are not an external
authentication service. Scheduled local alerts require a local sink; external
delivery is not automatically provided.

\section{Troubleshooting}

\begin{description}[style=nextline]
\item[No valid selection.] Inspect every raw repetition, status, feasibility,
fidelity and resource refusal. A missing/failed member cannot be dropped to
make a successful candidate. Preserve the failed study before changing a design.
\item[Resume refused.] Compare source, searcher state, declaration and budget
identities. Do not edit a checkpoint to force acceptance. Start an explicitly
new study when the protocol changes.
\item[Unsupported composition.] Check Table~\ref{tab:current-composition} and
the searcher's capabilities. Do not remove a declaration to disguise missing
method support.
\item[Search slower than grid.] Measure objective and orchestration/surrogate
cost separately. Tiny finite problems often justify enumeration; the worked
examples report this overhead instead of assuming HPO is faster.
\item[Model looks good but violates task constraints.] Check physical value,
derivative, boundary and identification conditions separately. Numerical
training loss alone cannot certify a scientific task.
\end{description}
